%% toc: P2 17 17.1,17.2,17.3,17.7
\section{First-order sequent adequacy}
\label{sec:p02-w02b-p02-019-lk-nd}

The comparison table \cref{tab:p02-fol-proof-systems} sets four first-order
calculi side by side.  For the sequent calculus two questions have to be kept
apart.  The first is which adequacy theorem holds, and under which language,
equality and set-theoretic hypotheses.  The second is how a derivation in the
sequent calculus is turned into a natural-deduction derivation and back, at
the same hypotheses.  This section answers both questions for the two
registered sequent calculi, the equality-free
\texttt{sheet:\allowbreak lk} and the identity calculus
\texttt{sheet:\allowbreak lk-eq}, against the natural-deduction
calculus \texttt{sheet:\allowbreak nd-fol}.  It enters the
proved adequacy theorems at their stated hypotheses, states the two-way
derivation maps for both hypotheses in one theorem, and writes out the one
case that the equality-free map had left as a one-sentence claim: the
natural-deduction rule \(\alpha\), which changes bound-variable names and has
no counterpart rule on the sequent sheet.  It also shows that the Axiom of
Choice is not needed for finite-sequent completeness of the equality-free
calculus, so that both sequent calculi have the same set-theoretic
hypothesis.

\emph{Inputs.}  A signature \(\Lang\) of any cardinality; a finite sequent
\(\Gamma\Rightarrow\Delta\) of finite multisets of \(\Lang\)-formulas; either
an \(\mathsf{LK}\)-derivation of it on one of the two sequent sheets or a
natural-deduction derivation of the translated formula.
\emph{Conclusions.}  The adequacy theorem that applies, with its hypotheses;
a derivation on the other calculus, obtained by a construction on finite
trees; and the limits of that construction.

\begin{hbdefinition}{def:p02-w02b-p02-019-hypotheses}{The two matched hypotheses}
A formula of \(\operatorname{Fm}(\Lang)\) is \emph{equality-free} when no
identity atom \(r=t\) occurs in it.  The two hypotheses are as follows.
\begin{enumerate}
 \item \emph{Equality-free hypothesis (EF).}  Every formula of the sequent is
   equality-free.  The sequent calculus is \texttt{sheet:lk}, whose language
   has no identity.  The natural-deduction calculus is
   \texttt{sheet:nd-fol} restricted to derivations all of whose formulas are
   equality-free; such a derivation never applies \(=I\) or \(=E\).
 \item \emph{Identity hypothesis (ID).}  Formulas may contain identity
   atoms.  The sequent calculus is \texttt{sheet:lk-eq}, written
   \(\mathsf{LK}_{=}\) as in \cref{def:p07-eqlk-domain}; the natural-deduction
   calculus is the whole of \texttt{sheet:nd-fol}, with \(=I\) and \(=E\).
\end{enumerate}
In both cases sequents are finite multisets on both sides, every term
instance is the total capture-avoiding substitution of
\cref{def:p02-fol-capture-avoiding-substitution}, and the eigenvariable
conditions are those of the sheets.  For a finite multiset \(\Delta\),
\(D_{\Delta}\) is the fixed disjunction of the distinct members of
\(\operatorname{supp}(\Delta)\) of \cref{def:p07-sequent-nd-translation},
with \(D_{\varnothing}=\bot\).  A sequent is \emph{valid} when every
structure with nonempty domain and every assignment that satisfy all members
of \(\operatorname{supp}(\Gamma)\) satisfy some member of
\(\operatorname{supp}(\Delta)\); identity is interpreted as actual identity.
\index{sequent calculus!matched hypotheses}
\end{hbdefinition}

The sequent sheets have no rule that renames bound variables, while the
natural-deduction sheet has the rule \(\alpha\): from \(\Gamma\vdash\varphi\)
infer \(\Gamma\vdash\psi\) whenever \(\varphi=_{\alpha}\psi\) in the sense of
\cref{def:p02-fol-alpha-equivalence}.  A translation of natural deduction
into the sequent calculus must therefore produce, for each alpha-equivalent
pair, sequent derivations of \(\varphi\Rightarrow\psi\) and
\(\psi\Rightarrow\varphi\), and then use them by Cut.  The next two lemmas
supply those derivations on \texttt{sheet:lk}.  The first records the
literal effect of substituting a variable that does not occur in a formula;
the second builds the derivations by induction on the definition of
\(=_{\alpha}\).

\begin{hblemma}{thm:p02-w02b-p02-019-renaming}{proved=annex:p02-w02b-p02-019-renaming}{Variable substitution identities}
Let \(\varphi\) be a formula and \(x\) a variable.
\begin{enumerate}
 \item \(\varphi[x/x]=\varphi\).
 \item If \(u,w\) are variables with \(u\notin\operatorname{Var}(\varphi)\)
   and \(w\notin\operatorname{Var}(\varphi)\cup\{u\}\), then
   \[
    \operatorname{Var}(\varphi[u/x])\subseteq\operatorname{Var}(\varphi)\cup\{u\}
    \qquad\text{and}\qquad
    (\varphi[u/x])[w/u]=\varphi[w/x].
   \]
\end{enumerate}
Each equation is literal identity of formulas, not alpha-equivalence.
Substitution of a variable keeps the relation symbol of every atom and
neither creates nor removes an identity atom, so \(\varphi[u/x]\) is equality-free exactly when
\(\varphi\) is.
\end{hblemma}

The hypothesis \(w\notin\operatorname{Var}(\varphi)\cup\{u\}\) is needed.
Substituting back the old variable can fail to restore the formula
literally: for \(\varphi=P(x)\land\forall x\,Q(x)\) and a variable \(u\)
distinct from \(x\), \(\varphi[u/x]=P(u)\land\forall x\,Q(x)\), and the
third clause of the substitution definition renames the binder in
\((\varphi[u/x])[x/u]\), which is therefore
\(P(x)\land\forall z\,Q(z)\) for the least admissible \(z\).  The lemma
avoids this by moving to a third variable \(w\).

\begin{hblemma}{thm:p02-w02b-p02-019-alpha-kernels}{proved=annex:p02-w02b-p02-019-alpha-kernels}{Equality-free alpha derivations}
Let \(\varphi=_{\alpha}\psi\).  Then \(\varphi\) is equality-free exactly
when \(\psi\) is.  If both are equality-free, \texttt{sheet:lk} derives
\[
 \varphi\Rightarrow\psi
 \qquad\text{and}\qquad
 \psi\Rightarrow\varphi
\]
by derivations that use no rule other than identity, weakening, the
propositional and quantifier rules, and Cut, and in which every formula is
equality-free.  Cut occurs only where two alpha-steps are composed; a single
generating step, and each step that applies one formula constructor to
already derived pairs, adds no Cut.
\end{hblemma}

\begin{hbtheorem}{thm:p02-w02b-p02-019-maps}{proved=annex:p02-w02b-p02-019-maps}{Two-way derivation maps}
Let \(\Gamma\Rightarrow\Delta\) be a finite sequent.
\begin{enumerate}
 \item Under the equality-free hypothesis, \texttt{sheet:lk} derives
   \(\Gamma\Rightarrow\Delta\) if and only if the equality-free restriction
   of \texttt{sheet:nd-fol} derives \(D_{\Delta}\) with all open assumptions
   in \(\operatorname{supp}(\Gamma)\).
 \item Under the identity hypothesis, \texttt{sheet:lk-eq} derives
   \(\Gamma\Rightarrow\Delta\) if and only if \texttt{sheet:nd-fol} derives
   \(D_{\Delta}\) with all open assumptions in
   \(\operatorname{supp}(\Gamma)\).
\end{enumerate}
In each direction and under each hypothesis the derivation on the other
calculus is built from the given one by recursion on the finite derivation
tree, one construction for each rule of the sheet, including Mix, the
eigenvariable rules, the \(\alpha\) rule and the nonvacuous-discharge
requirement.  Under the equality-free hypothesis no constructed derivation
contains an identity atom.  The maps establish equivalence of derivability;
they do not identify proof trees, and the natural-deduction-to-sequent
direction uses Cut.
\end{hbtheorem}

Clause~2 is \cref{thm:p07-eqlk-nd}.  Clause~1 restates
\cref{thm:p07-lk-nd-equivalence}; its proof in the annex obtains the
equality-free maps by restricting the identity constructions and replacing
the alpha case by \cref{thm:p02-w02b-p02-019-alpha-kernels}.  This also
covers the Mix rule of \texttt{sheet:lk}, which the earlier equality-free
proof does not list among its cases.

\begin{hbcorollary}{thm:p02-w02b-p02-019-choice-free}{proved=annex:p02-w02b-p02-019-choice-free}{Choice-free adequacy of equality-free LK}
Let \(\Lang\) be a signature of any cardinality and let
\(\Gamma\Rightarrow\Delta\) be a finite sequent of equality-free
\(\Lang\)-formulas.  Then \texttt{sheet:lk} derives
\(\Gamma\Rightarrow\Delta\) if and only if the sequent is valid.  The proof
uses no form of the Axiom of Choice.  It uses Cut, and no cut-free
derivation is claimed.
\end{hbcorollary}

The proved sequent-calculus adequacy results and their hypotheses are
collected in \cref{tab:p02-w02b-p02-019-lk-adequacy}.  The inherited
\cref{thm:fol-lk-completeness} keeps its status as a proof sketch and its
stated assumption of the Axiom of Choice; the table lists the proved
theorems that cover its content.

\begingroup\small
\begin{longtable}{@{}p{.17\textwidth}p{.37\textwidth}p{.38\textwidth}@{}}
\caption{Adequacy of the two sequent calculi at their hypotheses.}
\label{tab:p02-w02b-p02-019-lk-adequacy}\\
\toprule
Property & Equality-free: \texttt{sheet:lk} & Identity: \texttt{sheet:lk-eq} \\
\midrule
\endfirsthead
\caption[]{Adequacy of the two sequent calculi at their hypotheses (continued)}\\
\toprule
Property & Equality-free: \texttt{sheet:lk} & Identity: \texttt{sheet:lk-eq} \\
\midrule
\endhead
\midrule\multicolumn{3}{r}{\emph{continued on next page}}\\
\endfoot
\endlastfoot
Language & No identity atoms; any signature. & Primitive identity; any
signature. \\
Soundness & Proved, \cref{thm:fol-lk-soundness}; a Mix step is a
Cut after contractions. & Proved for every rule, including Mix,
\cref{thm:p07-eqlk-soundness}. \\
Finite-sequent completeness & Proved with Cut,
\cref{thm:p07-canonical-lk-completeness}: Choice assumed for an arbitrary
signature, not for a finite or countable one.  Proved without Choice for
every signature, \cref{thm:p02-w02b-p02-019-choice-free}. & Proved with Cut
for every signature and without Choice,
\cref{thm:p07-eqlk-completeness}, by reduction to the finite signature of the
sequent. \\
Route of the completeness proof & Natural-deduction completeness for the
equality-free language, then the translation of
\cref{thm:p07-lk-nd-equivalence}; the finite-signature reduction of
\cref{thm:p02-w02b-p02-019-choice-free} removes Choice. & Natural-deduction completeness
with identity for countable signatures,
\cref{cor:fol-completeness-countable}, then
\cref{thm:p07-eqlk-nd}. \\
Two-way natural-deduction maps & \cref{thm:p02-w02b-p02-019-maps}(1),
\cref{thm:p07-lk-nd-equivalence}; alpha case by
\cref{thm:p02-w02b-p02-019-alpha-kernels}. &
\cref{thm:p02-w02b-p02-019-maps}(2), \cref{thm:p07-eqlk-nd}; alpha case by
\cref{thm:p07-eqlk-alpha}. \\
Cut & No cut-elimination theorem is proved for this presentation; the
cut-elimination theorem \cref{thm:p07-classical-cut-elimination} concerns the
defined-negation calculus \texttt{sheet:lk-defneg}. & Cut and Mix are not
admissible: the Cut/Mix-free fragment does not derive \(a=b\Rightarrow b=a\),
\cref{thm:p07-eqlk-cut-failure}. \\
Inherited sketch & \cref{thm:fol-lk-completeness}, sketch, assumes Choice. &
The same sketch includes equality; its equality translation is supplied by
\cref{thm:p07-eqlk-nd}. \\
\bottomrule
\end{longtable}
\endgroup

\begin{hbexample}{def:p02-w02b-p02-019-example}{One alpha step in both directions}
Let \(P\) be a unary relation symbol and \(x,y\) distinct variables.  The
natural-deduction derivation
\[
 \forall x\,P(x)\vdash\forall x\,P(x)\ \text{(assumption)},\qquad
 \forall x\,P(x)\vdash\forall y\,P(y)\ (\alpha)
\]
is correct because \(\forall x\,P(x)=_{\alpha}\forall y(P(x)[y/x])\) is a
generating pair, with \(y\notin\operatorname{Var}(P(x))=\{x\}\), and
\(P(x)[y/x]=P(y)\).

\emph{Natural deduction to sequents.}  The assumption becomes the identity
sequent \(\forall x\,P(x)\Rightarrow\forall x\,P(x)\).  For the \(\alpha\)
step, choose a variable \(z\) distinct from \(x\) and \(y\).  The kernel of
\cref{thm:p02-w02b-p02-019-alpha-kernels} for this generating pair is
\[
\begin{array}{rll}
 1.&P(z)\Rightarrow P(z)&\text{identity}\\
 2.&\forall x\,P(x)\Rightarrow P(z)&\forall L\ 1\\
 3.&\forall x\,P(x)\Rightarrow\forall y\,P(y)&\forall R\ 2
\end{array}
\]
The \(\forall L\) step uses the term \(z\), since \(P(x)[z/x]=P(z)\); the
\(\forall R\) step uses the eigenvariable \(z\), since \(P(y)[z/y]=P(z)\),
and \(z\) is free in neither \(\forall x\,P(x)\) nor \(\forall y\,P(y)\).  A
Cut of the identity sequent against this kernel on the formula
\(\forall x\,P(x)\) gives \(\forall x\,P(x)\Rightarrow\forall y\,P(y)\).
Here \(D_{\{\forall y P(y)\}}=\forall y\,P(y)\), so the closing step of the
map is a Cut with an identity sequent.

\emph{Sequents to natural deduction.}  Applied to the kernel, the map of
\cref{thm:p02-w02b-p02-019-maps}(1) sends the identity leaf to the
assumption \(P(z)\) and \(\forall L\) to \(\forall E\) with the term \(z\),
giving a derivation of \(P(z)\) from \(\forall x\,P(x)\).  At \(\forall R\)
the retained succedent is empty, so \(D=\bot\).  The construction renames
the eigenvariable to a variable \(z'\) absent from the derivation so far and
from the sequent, and then produces the following derivation, with open
assumption \(\forall x\,P(x)\):
\[
\begin{array}{rll}
 1.&\forall x\,P(x)&\text{assumption}\\
 2.&P(z')&\forall E\ 1\\
 3.&\bot\lor P(z')&\lor I_2\ 2\\
 4.&\neg(\bot\lor\forall y\,P(y))&\text{assumption, label }a\\
 5.&\bot&\text{assumption, label }b\\
 6.&\bot\lor\forall y\,P(y)&\lor I_1\ 5\\
 7.&\bot&\neg E\ 6,4\\
 8.&\neg\bot&\neg I\ 7\text{, discharges }b\\
 9.&\bot&\text{assumption, label }c\\
 10.&\bot&\neg E\ 9,8\\
 11.&P(z')&\bot E\ 10\\
 12.&P(z')&\text{assumption, label }d\\
 13.&P(z')&\lor E\ 3,\ 9\text{--}11,\ 12\text{, discharges }c,d\\
 14.&\forall z'\,P(z')&\forall I\ 13\\
 15.&\forall y\,P(y)&\alpha\ 14\\
 16.&\bot\lor\forall y\,P(y)&\lor I_2\ 15\\
 17.&\bot&\neg E\ 16,4\\
 18.&\bot\lor\forall y\,P(y)&\text{RAA }17\text{, discharges }a\\
 19.&\bot&\text{assumption, label }e\\
 20.&\forall y\,P(y)&\bot E\ 19\\
 21.&\forall y\,P(y)&\text{assumption, label }f\\
 22.&\forall y\,P(y)&\lor E\ 18,\ 19\text{--}20,\ 21\text{, discharges }e,f
\end{array}
\]
At line~14 the open assumptions are lines~1 and~4, in neither of which
\(z'\) is free.  Line~15 is a generating pair because
\(y\notin\operatorname{Var}(P(z'))\).  Every discharge marks an assumption
that is used above it.  Lines~19--22 convert \(\bot\lor\forall y\,P(y)\) to
\(D_{\{\forall y P(y)\}}\).  Translating the Cut grafts the identity
derivation of \(\forall x\,P(x)\) into line~1, which leaves the derivation
unchanged.

The round trip turns a two-node derivation into one with twenty-two lines.
The maps preserve derivability, not the shape or size of derivations.
\end{hbexample}

\begin{hbwarning}{def:p02-w02b-p02-019-limits}{What the matched maps do not give}
\begin{enumerate}
 \item The natural-deduction-to-sequent direction uses Cut: at every
   elimination rule, at \(\alpha\), and in the last step that turns
   \(\operatorname{supp}(\Gamma)\Rightarrow D_{\Delta}\) into
   \(\Gamma\Rightarrow\Delta\).  Under the identity hypothesis Cut cannot be
   removed: \(a=b\Rightarrow b=a\) is valid and derivable, but not without Cut
   or Mix (\cref{thm:p07-eqlk-cut-failure}).  The adequacy theorems in
   \cref{tab:p02-w02b-p02-019-lk-adequacy} are therefore statements about the
   calculi with Cut.
 \item The translation \(\Gamma\Rightarrow\Delta\mapsto
   (\operatorname{supp}(\Gamma),D_{\Delta})\) forgets multiplicities and the
   order of succedent formulas.  The maps are not inverse to each other and
   do not preserve height, size, or normality.
 \item The equality-free and identity maps are not interchangeable.  An
   equality-free sequent derivable on \texttt{sheet:lk-eq} need not come with
   an equality-free derivation from clause~2, since its natural-deduction
   derivation may use \(=I\) and \(=E\).  Clause~1 and
   \cref{thm:p02-w02b-p02-019-choice-free} give an equality-free derivation
   directly.
 \item The adequacy results are for finite sequents.  A consequence from an
   infinite premise set is reached through natural-deduction completeness,
   \cref{thm:fol-completeness}, which assumes Choice for an arbitrary
   signature, or \cref{cor:fol-completeness-countable} for a countable one.
\end{enumerate}
\end{hbwarning}

\emph{Relationships and sources.}  The Hilbert and tableau rows of \cref{tab:p02-fol-proof-systems} have their
own adequacy theorems, \cref{thm:p02-hfa-adequacy} and
\cref{thm:p02-fol-tableau-adequacy}.  The comparison of a sequent with a
formula follows Gentzen, who calls a sequent equivalent to the implication
from the conjunction of its antecedent formulas to the disjunction of its
succedent formulas, with separate clauses for empty sides
\cite[Abschnitt~V, \S~1, pp.~417--418]{Gentzen1935b}.  Gentzen's own
equivalence proof for the classical calculi does not map sequent derivations
directly into natural deduction: it translates natural deduction into the
sequent calculus, the sequent calculus into a Hilbert-type calculus, and that
calculus back into natural deduction
\cite[Abschnitt~V, \S~6, p.~429]{Gentzen1935b}.  The direct sequent-to-natural
deduction map of \cref{thm:p02-w02b-p02-019-maps} is the construction of
\cref{thm:p07-lk-nd-equivalence,thm:p07-eqlk-nd}.
